\subsubsection{2026 年 7 月 29 日至 8 月 25 日量化训练}

\begin{longtable}{p{1.4cm}p{1.55cm}p{3.1cm}p{4.6cm}p{3.4cm}}
\caption{四周训练计划与验收标准}\label{tab:four-weeks}\\
\toprule
时间 & 负责人 & 训练任务 & 产物 & 验收标准\\
\midrule
\endfirsthead
\toprule
时间 & 负责人 & 训练任务 & 产物 & 验收标准\\
\midrule
\endhead
7/29--8/4 & 队长 & 完成 2022A 动力学最小复现；比较 Euler、RK4 与自适应积分。 & 方程说明、两套代码、步长收敛图、守恒与回代误差表。 & 小样例误差随步长按预期下降；关键结果可复算；接口测试全部通过。\\
7/29--8/4 & 组员1 & 完成 2024B 的枚举基线与 0--1/MILP 表达。 & 决策树、模型文件、枚举核验、不可行情形清单。 & 小规模最优值一致；每条约束有题意来源；报告最优界或间隙。\\
7/29--8/4 & 组员2 & 完成 2022C 数据审计、可视化与分类基线。 & 数据字典、清洗审计、pipeline、验证报告。 & 预处理只在训练集拟合；异常处理可追溯；至少三类基线同口径比较。\\
8/5--8/11 & 队长 & 完成热/扩散方程 FDM 与复杂几何 FEM 对照。 & 网格文件、边界测试、网格收敛表。 & 至少三档网格；边界和量纲单测通过；说明 FDM/FEM 选型边界。\\
8/5--8/11 & 组员1 & 完成 2023B 图/几何基线及 GA、PSO 公平对照。 & 基线求解器、随机算法报告、可行性修复测试。 & 不少于 20 个种子；预算相同；报告中位数、离散度、时间和可行率。\\
8/5--8/11 & 组员2 & 完成 2023C 时序滚动验证和补货特征接口。 & 持续性/季节基线、ARIMA、树模型、滚动评估。 & 无未来信息泄漏；至少四个误差指标；失败日期有原因分析。\\
8/12--8/18 & 全队 & 选择一题进行 36 小时压缩演练，贯通模型、代码、图表、论文和附件。 & 可编译论文、运行入口、结果清单、复盘表。 & 所有正文数字可定位到输出；无未定义引用；第三人按说明复现关键表。\\
8/19--8/25 & 全队 & 进行 72 小时 A/B/C 三选一全流程模拟并交叉审稿。 & 最终论文、源代码、附件、审计 JSON、个人复盘。 & 12 小时内定题；24 小时出基线；48 小时主结果冻结；终稿通过结构、引用、检验和复现门。\\
\bottomrule
\end{longtable}

